\section{补充数据}\label{sec:appendix-data}

本附录提供正文中未详细列出的补充数据与推导过程，供读者参考。

\subsection{数据处理公式推导}

正文中使用的绩点转换公式，其分段函数形式如下：
\begin{equation}\label{eq:gpa-conversion}
  G(s) =
  \begin{cases}
    4.0, & s \geq 90, \\
    3.0, & 80 \leq s < 90, \\
    2.0, & 70 \leq s < 80, \\
    1.0, & 60 \leq s < 70, \\
    0.0, & s < 60.
  \end{cases}
\end{equation}
其中 $s$ 为百分制成绩，$G(s)$ 为对应绩点。

\subsection{补充数据表}

\cref{tab:appendix-sample} 列出了更多测试样本的详细数据。

\begin{table}[!htb]
  \caption{补充测试数据}
  \label{tab:appendix-sample}
  \centering
  \begin{tabular}{@{}cccc@{}} \toprule
    \textbf{样本编号} & \textbf{输入值} & \textbf{输出值} & \textbf{误差} \\ \midrule
    1 & 0.50 & 0.48 & 0.02 \\
    2 & 1.00 & 0.97 & 0.03 \\
    3 & 1.50 & 1.52 & 0.02 \\
    4 & 2.00 & 1.98 & 0.02 \\
    5 & 2.50 & 2.53 & 0.03 \\ \bottomrule
  \end{tabular}
\end{table}
